\documentclass[11pt,a4paper]{article}
\usepackage[margin=23mm]{geometry}
\usepackage{amsmath,amssymb,booktabs,array,longtable,hyperref}
\hypersetup{hidelinks}
\emergencystretch=3em
\newcommand{\Lc}{\textbf{(c)}\ }
\newcommand{\Ld}{\textbf{(d)}\ }
\newcommand{\dd}{\mathrm d}
\newcommand{\scan}{\texttt{onebody\_\allowbreak prec\_\allowbreak scan.py}}
\newcommand{\edge}{\texttt{onebody\_\allowbreak prec\_\allowbreak edge.py}}
\newcommand{\ids}{\texttt{onebody\_\allowbreak prec\_\allowbreak identities.py}}
\newcommand{\ana}{\texttt{onebody\_\allowbreak prec\_\allowbreak analysis.py}}
\title{The Toda one-body coefficient at higher precision:\\
independent finite-sea extrapolation and an analytic edge reduction}
\author{Research note for Gustav Delius}
\date{4 November 2026}
\begin{document}
\maketitle
\begin{abstract}
\Lc\ An unconstrained extrapolation of new finite-sea Bethe data gives
$C=0.1591549414\pm 8.23\times10^{-6}$.
The mesh/roundoff, $h\to0$, and final amplitude-tail components are $5.68\times10^{-8}$, $7.70\times10^{-6}$, and $4.73\times10^{-7}$, respectively. This reduces the previous $0.00334$ envelope by a factor of 406.
The uncertainties are deterministic numerical diagnostics, added linearly,
not statistical confidence intervals or rigorous bounds.
\Lc\ A separately checked half-line calculation derives $1/(2\pi)$ from
the momentum logarithm and the harmonic Toda vacuum sound correction.
The numerical fit does not impose this constant.
\Ld\ The result concerns the stipulated Toda edge limit and inherited
classical Boussinesq units. A uniform interchange with the fixed-nonzero-$P$
semiclassical limit, and identification with a completed quantum bad-Boussinesq
field theory, remain additional questions.
\end{abstract}

\section{Scope, normalization, and checked inputs}
Labels mean (a) proved in the literature, (b) literature heuristic,
(c) derived or checked in this project, and (d) unestablished identification
or guess. All numerical (c) statements mean evidence at the stated precision,
not existence theorems. The complete internal notes
\texttt{onebody\_dispersion\_test.tex} and \texttt{frame\_dictionary.tex},
their specified \texttt{onebody\_*} and \texttt{frame\_*} implementations,
and the relevant classical Toda and earlier half-line sections were read.
No external reference, DOI, arXiv entry, full text, or abstract is cited or
claimed as independently verified here. No older script is executed: several
write files outside this task's permitted prefixes.

\paragraph{Inherited units \Lc\ (\edge, \ana).}
Use $a=1$, $b=e^{-1/2}$, $h=\hbar_T=3b\hbar_B$ and
\begin{equation}
 P_B=\frac{\overline P}{6\sqrt3 b},\qquad
 E_B=\frac{\overline E}{6\sqrt3 b^2},\qquad
 C_h(y)=3\frac{\overline E(P_0(\kappa))-E_0(\kappa)}{hP_0(\kappa)},
 \quad \kappa=\sqrt{hy/b}.
 \label{eq:quantity}
\end{equation}
Here
\begin{equation}
 P_0=4b(\kappa\cosh\kappa-\sinh\kappa),\qquad
 E_0=2b^2(\sinh\kappa\cosh\kappa-\kappa),\qquad
 \frac{\dd E_0}{\dd P_0}=b\frac{\sinh\kappa}{\kappa}.
\end{equation}
The reported coefficient is $C=\lim_{y\to\infty}\lim_{h\to0}C_h(y)$.
It multiplies $\hbar_B|P_B|$ in energy, or is order $\hbar_B^0$ in
$E_B/\hbar_B$. Parity supplies the opposite edge.
The inherited continuum equation is the \emph{bad} one,
$u_{tt}=u_{xx}+u_{xxxx}+3(u^2)_{xx}$, with
$u=2(\log\tau)_{xx}$, $k=\sinh\beta$, $v=\epsilon\cosh\beta$ and
$\kappa=\sqrt3\beta$ to leading order. There is no change of regulator,
time normalization, sound-line convention, or quantum charge identification.

\section{Fresh finite-sea calculation}
\paragraph{Equations \Lc\ (\scan, \ids).}
With $t=2\pi h$, the sea operator and its two solutions are
\begin{align}
 \Theta(d)&=2h\operatorname{Im}\log\Gamma(1+id/h)+2d\log h,
 &K&=\Theta',\\
 (\mathcal Af)(x)&=tf(x)-\int_{-q}^q K(x-z)f(z)\,\dd z,
 &\mathcal A\rho&=1,\quad \mathcal A\eta=x,\quad\int\rho=1.
\end{align}
Define $D=1+\int K(p-x)\rho(x)\dd x$,
$V=p+\int K(p-x)\eta(x)\dd x$, and $v_s=V(q)/D(q)$.
The edge-relative observables are
\begin{align}
 \overline P(p)&=p-q+\int[\Theta(p-x)-\Theta(q-x)]\rho(x)\,\dd x,\\
 \overline E(p)&=\tfrac12(p^2-q^2)
       +\int[\Theta(p-x)-\Theta(q-x)]\eta(x)\,\dd x.
\end{align}
We solve $\overline P(p)=P_0(\kappa)$ before evaluating each energy.
The direct group-velocity coefficient is
\begin{equation}
 C_{v,h}(y)=\frac3h\left[\frac{V(p)}{D(p)}
                         -b\frac{\sinh\kappa}{\kappa}\right].
\end{equation}

\paragraph{Subtractions and conditioning \Lc.}
The full classical lattice energy, not merely its cubic approximation,
is subtracted. To reduce cancellation, the code first forms
\begin{equation}
 \overline E-b\overline P
 =(p-q)\left(\frac{p+q}{2}-b\right)
 +\int[\Theta(p-x)-\Theta(q-x)](\eta-b\rho)(x)\,\dd x,
 \label{eq:stable}
\end{equation}
and subtracts $E_0-bP_0$, calculated at 60 decimal digits.
This is algebraically the original energy comparison: the quantum sound
shift has \emph{not} been removed. Direct evaluation of the unsubtracted
formula is checked too. The dense solves themselves use IEEE double
precision; high-precision fits do not upgrade those input data.

\paragraph{Meshes and extrapolation \Lc\ (\scan, \ana).}
The eleven $h$ values are
$.08,.06,.04,.03,.02,.015,.01,.0075,.005,.00375,.0025$.
The base mesh is $N=\lceil10/h+200\rceil$; refinements increase this by
25\% at $.01,.005,.0025$, reaching $N=5250$.
Fifteen momentum labels cover $y=.25$ through $128$.
All sea solves check normalization, $P(q)=\pi h$, monotonic momentum,
and the fixed-$P$ inversion. The exact derivative identity below is also
checked independently at 60 digits, by complex trigamma differentiation and
separate finite differences in $q$, not by substituting the identity into itself.

The central $h$ intercept is a quartic least-squares fit to the last seven
$h$ values. Alternatives are a cubic through the last six, quintic through
the last eight, quadratic through the last four, and cubic plus either
$h\log h$ or $h^2\log h$ through the last seven. All coordinates are scaled
before fitting and all small linear systems use 60-digit arithmetic.
The $h$ error is the maximum intercept change. Lower-order Richardson-type
fits are retained as conservative truncation diagnostics.
The mesh diagnostic is the largest refined-minus-base difference plus
$2\,10^{-10}$ for roundoff, multiplied by the largest $\ell^1$ norm of the
$h$-extrapolation weights. These errors are subsequently propagated through
the amplitude-fit weights.

\begin{table}[ht]\centering\small
\begin{tabular}{rrrrr}\toprule
$y$ & $C_E(y)$ & mesh/roundoff & $h$ extrapolation & half-line difference\\\midrule
4 & 0.16093283146 & $9.8\times10^{-9}$ & $2.0\times10^{-9}$ & $2.7\times10^{-10}$ \\
8 & 0.16002007918 & $9.4\times10^{-9}$ & $1.3\times10^{-8}$ & $6.4\times10^{-11}$ \\
16 & 0.15958191769 & $1.1\times10^{-8}$ & $8.9\times10^{-8}$ & $1.4\times10^{-10}$ \\
32 & 0.15936706868 & $1.3\times10^{-8}$ & $6.2\times10^{-7}$ & $3.2\times10^{-9}$ \\
64 & 0.15926060097 & $1.9\times10^{-8}$ & $4.4\times10^{-6}$ & $7.2\times10^{-8}$ \\
\bottomrule\end{tabular}
\caption{Fresh finite-sea energy intercepts (c). The last column compares with independently inverted half-line equations, not a fitted target constant. Both energy and velocity differences through $y=64$ are inside the separately assigned mesh-plus-$h$ diagnostics.}
\end{table}

\section{Analytic edge reduction and removal of the logarithm}
\paragraph{An exact thermodynamic identity \Lc\ (\ids, \scan).}
For this paragraph $q$ varies at fixed $h$ and $a$ is determined by counting;
it is not held fixed when differentiating $q$. Differentiating the sea
equations and integrating by parts gives
\begin{equation}
 \mathcal A\eta'=1-\eta(q)[K(x-q)+K(x+q)],\qquad
 \mathcal A\partial_q\rho=\rho(q)[K(x-q)+K(x+q)].
\end{equation}
Consequently $\eta'=\rho-v_s\partial_q\rho$, and its exterior version is
\begin{equation}
 \boxed{\partial_pV=D-v_s\partial_qD.} \label{eq:identity}
\end{equation}
Independent $N=32,48$ mpmath quadratures at $h=.7,q=1$, with a centered
$q$ difference of step $10^{-15}$, give maximum residuals
$3.21\,10^{-19}$ and $1.13\,10^{-28}$, respectively.

\paragraph{Edge shape and its normalization \Lc\ (\edge).}
Put $u=(p-q)/h$ and $n=1/a$. Classical matching gives
$n=-1/[2\log(q/2)]$ and, to leading edge order,
\begin{equation}
 D(p,q)=A(q,h)f(u)+o(A),\qquad
 A=n\sqrt{2h/q},\qquad F(u)=\int_0^u f(w)\dd w=\tfrac43 y(u)^{3/2}.
 \label{eq:edgeform}
\end{equation}
Then $\overline P=hAF$ and $v_{s,0}=q/(2n)$.
The logarithmic derivative of $A$ satisfies
$1-v_{s,0}\partial_q\log A=1/(4n)$, checked symbolically.
In \eqref{eq:identity} the $-\partial_u/h$ part of $\partial_q$
cancels $\partial_pD$. Matching $V(q)=v_sD(q)$ fixes the integration constant:
\begin{equation}
 V-v_sD=\frac{hA}{4n}F+o(hA),\qquad
 \frac VD=v_s+\frac{h}{4n}\frac Ff+o(h).
 \label{eq:velocityedge}
\end{equation}
These are asymptotic statements in the matched edge expansion, rather than
an assertion of uniformity in arbitrarily large $y$ at finite $h$.

\paragraph{Correct half-line factor \Lc\ (\edge).}
The integral representation of the digamma kernel gives, for $\omega\ne0$,
\begin{equation}
 \widehat K_0(\omega)=-\frac{2\pi}{e^{|\omega|}-1},\qquad
 \Sigma(\omega)=\frac{2\pi}{1-e^{-|\omega|}}
              =2\pi g(i\omega)g(-i\omega),
\end{equation}
where
\begin{equation}
 g(s)=s^{-1/2}\Gamma\left(1+\frac{s}{2\pi}\right)
 \exp\left[-\frac{s}{2\pi}\log\frac{s}{2\pi e}\right],\qquad \Re s>0.
\end{equation}
The kernel transform follows by transforming the cosine in
$K_0(x)=-2\gamma+2\int_0^\infty(1-\cos xt)/(e^t-1)\dd t$;
the additive zero-frequency distribution is fixed by bulk matching.
The factor product is independently checked at five nonzero frequencies
to $10^{-64}$. This corrects the Fourier scale in the older Phase~2i note;
that note's printed transform is not used.

The half-line solution has the interior transform proportional to the
inverse interior factor divided by the zero-frequency pole. Multiplication
by $\Sigma$ leaves the exterior factor divided by that pole. Thus
$\mathcal L[f]\propto g(s)/s$, and integration fixes
\begin{equation}
 \boxed{\mathcal L[F](s)=\sqrt\pi\,\frac{g(s)}{s^2},\qquad
 \mathcal L[J](s)=\sqrt\pi\,\frac{g(s)}{s^3},\quad
 J(u)=\int_0^u F(w)\dd w.} \label{eq:laplace}
\end{equation}
The normalization follows from $F\sim4u^{3/2}/3$, whose transform is
$\sqrt\pi s^{-5/2}$. In particular, omitting $\sqrt\pi$ fails the classical
normalization. The branch is the principal logarithm in the right half-plane.
Independent de Hoog (50 digits) and Talbot (70 digits) inversions agree
within $3\,10^{-47}$ in all reported quantities at
$y=4,8,16,32,64,128,256,512,1024$.
The fresh finite-sea intercepts test both the momentum shape and the energy
consequences of this half-line construction.

\paragraph{Vacuum sound coefficient \Lc\ (\edge, \scan).}
Expanding the inherited Toda Hamiltonian
$H=\sum_n[p_n^2/2+e^{-(x_{n+1}-x_n)}]$ about spacing $a$ gives
$\omega(k)=2b|\sin(k/2)|$. Its harmonic zero-point energy per particle is
\begin{equation}
 e(a,h)=b^2+h\frac1{4\pi}\int_0^{2\pi}\omega(k)\dd k+O(h^2)
       =b^2+\frac{2hb}{\pi}+O(h^2).
\end{equation}
The long-wave elastic equation with unit particle mass gives
$v_s=a\sqrt{\partial_a^2e}=ab+ah/(4\pi)+O(h^2)$.
The integral and derivative are checked symbolically. This supplies the
vacuum normalization independently of the proposed soliton coefficient;
the finite-sea sound extrapolation tests it separately.
No arbitrary multiple of the sound line is adjusted to obtain agreement.

\paragraph{Energy and leading logarithm \Lc\ (\edge).}
At $a=1$, integrating \eqref{eq:velocityedge} at fixed dressed momentum yields
\begin{equation}
 C_E(y)=\frac3{4\pi}+\frac{3J(u(y))}{4F(u(y))}-\frac{3y}{10},\qquad
 C_v(y)=\frac3{4\pi}+\frac y2[u'(y)-1].
 \label{eq:edgeenergy}
\end{equation}
The $3y/10$ is exactly the leading edge limit of the classical nonlinear
energy already removed in \eqref{eq:quantity}.
Expanding the independently factored transform at small $s$ gives
\begin{align}
 g(s)&=s^{-1/2}\left[1+\frac{s}{2\pi}
 (1+\log(2\pi)-\gamma-\log s)+O(s^2\log^2s)\right],\\
 u(y)&=y-A_0\log y+c_0+\frac{d_0}{y}+O(y^{-2}\log^m y),\label{eq:u}\\
 A_0&=\frac1{2\pi},\qquad
 c_0=\frac{1-\log(8\pi)}{2\pi},\qquad
 d_0=\frac1{48}-\frac1{4\pi^2}.
\end{align}
The transform inversion uses
$\mathcal L^{-1}[s^{-\alpha}\log^j s]
=(-\partial_\alpha)^j[u^{\alpha-1}/\Gamma(\alpha)]$.
SymPy differentiates these expressions and checks the inverse-series
coefficients, including cancellation of $(\log y)/y$ in $u(y)$.

In the energy integral use $\dd u=u'(y)\dd y$ \emph{before} extrapolation:
\begin{equation}
 J=\frac8{15}y^{5/2}-\frac{8A_0}{9}y^{3/2}
       -\frac{8d_0}{3}y^{1/2}+\text{constant}+O(y^{-1/2}\log^m y).
\end{equation}
Thus both the leading classical term and the known logarithmic displacement
are removed analytically; no unknown leading $\log y$ energy coefficient is fitted.
The constant of integration is allowed as a $y^{-3/2}$ term in numerical fits.
Equations \eqref{eq:edgeenergy}--\eqref{eq:u} give
\begin{align}
 C_E(y)&=\frac3{4\pi}-\frac{A_0}{2}-\frac{3d_0}{2y}
                  +O(y^{-3/2}),\\
 C_v(y)&=\frac3{4\pi}-\frac{A_0}{2}-\frac{d_0}{2y}+o(y^{-1}),\\
 \boxed{C=\frac3{4\pi}-\frac1{4\pi}=\frac1{2\pi}.} \label{eq:answer}
\end{align}
This is a project analytic derivation within the matched Bethe edge problem,
with the harmonic vacuum normalization stated explicitly. It is not an
assumed fit intercept or a new literature theorem.

\section{Numerical verdict and separated uncertainty}
\paragraph{Free intercept test \Lc\ (\ana).}
The finite-sea energy fits first subtract the derived tail $-3d_0/(2y)$.
The fitted constant remains free. We vary $y$ windows $[4,32]$, $[8,32]$,
$[8,64]$, $[16,64]$ and residual bases
\[
 \{1,y^{-2}\},\quad\{1,y^{-2},y^{-3}\},\quad
 \{1,y^{-3/2},y^{-2}\},\quad\{1,y^{-2},y^{-2}\log y\}.
\]
The central fit uses $[8,32]$ and the second basis.
The amplitude error is the maximum intercept change among all sixteen fits.
The mesh and $h$ components are separately propagated with absolute fit
weights, taking the largest result over those fits, and the three maxima
are added. This deliberately conservative total also retains lower-order
$h$ fits and a possible integration-constant tail.
Here ``amplitude error'' denotes the final large-$y$ tail of the edge
coefficient; the actual lattice amplitude $\kappa^2=hy/b$ has already
been extrapolated to zero in the first stage.

\Lc\ The primary result and its direct velocity cross-check are
\begin{align*}
 C_E&=0.15915494136\pm 8.23\times10^{-6},\\
 C_v&=0.15915490084\pm 4.62\times10^{-5}.
\end{align*}
\begin{center}
\begin{tabular}{lrr}\toprule
Error component & energy & velocity\\\midrule
Mesh and roundoff & $5.68\times10^{-8}$ & $7.48\times10^{-8}$ \\
$h\to0$ extrapolation & $7.70\times10^{-6}$ & $4.49\times10^{-5}$ \\
Amplitude-tail extrapolation & $4.73\times10^{-7}$ & $1.27\times10^{-6}$ \\
Linear total & $8.23\times10^{-6}$ & $4.62\times10^{-5}$ \\
\bottomrule\end{tabular}\end{center}
The independent sound intercept is $3(v_s-b)/h=0.238732414641$, with $h$-fit spread $2.24\times10^{-10}$; its analytic comparator is $3/(4\pi)$.
\begin{center}\small
\begin{tabular}{lrrl}\toprule
Hypothesis & value & discrepancy/envelope & verdict\\\midrule
$1/(2\pi)$ & 0.1591549431 & 0.000211 & compatible \\
$1/6$ & 0.1666666667 & 912 & excluded \\
$\sqrt3/(2\pi)^2$ & 0.0438733595 & $1.40\times10^4$ & excluded \\
$(1+10^{-3})/(2\pi)$ & 0.1593140980 & 19.3 & excluded \\
$(1+10^{-4})/(2\pi)$ & 0.1591708586 & 1.93 & excluded \\
$(1+10^{-5})/(2\pi)$ & 0.1591565346 & 0.194 & compatible \\
\bottomrule\end{tabular}\end{center}
For $C=(1+\delta)/(2\pi)$ the finite-sea envelope gives $-5.174\times10^{-5}\leq\delta\leq 5.172\times10^{-5}$. Thus arbitrarily small nonzero offsets remain numerically indistinguishable. The alternative $\sqrt3/(2\pi)^2$ is about $0.04387$, so it is not actually a nearby value.
The twelve free-intercept high-precision half-line fits give the range $[0.15915494308687,0.15915494310110]$, including variation of the fitting window and a $y^{-3/2}$ tail. This is an additional asymptotic diagnostic, separate from the primary finite-sea budget.


\paragraph{Analytic versus numerical precision.}
\Lc\ The half-line inversion has vastly smaller arithmetic error than the
finite-sea extrapolation. Fitting its high-precision values with a free
intercept provides a useful additional diagnostic, but does not justify
replacing the finite-sea error budget by a 50-digit error bar.
Agreement with \eqref{eq:answer} is compatibility within a deterministic
envelope, not a probability that two real numbers are exactly equal.
For the model $C=(1+\delta)/(2\pi)$ the interval on $\delta$ is explicitly
reported; ``small'' without a specified magnitude cannot be rejected.

\paragraph{Unbounded qualifications \Ld.}
The derivation uses smooth matched edge asymptotics and their differentiated
form. The numerical errors do not bound failure of that asymptotic ansatz,
or an exchange of $h\to0$, $y\to\infty$ with the fixed-nonzero-$P$
semiclassical limit. In particular $y\to\infty$ at fixed $h$ exits the
small-amplitude region. Neither the unresolved field determinant nor a
quantum correction to the Toda--Boussinesq identification has an assigned
error bar here. These are distinct from the numerical coefficient of the
explicitly defined nested Bethe limit.

\section{Checks, failures recorded, and reproduction}
\begin{center}\small
\begin{tabular}{>{\raggedright\arraybackslash}p{.45\linewidth}c>{\raggedright\arraybackslash}p{.36\linewidth}}
\toprule Claim & Label & Check\\\midrule
Fixed-$P$ energy and direct velocity & (c) & \scan; refined meshes\\
Exact $V'=D-v_s\partial_qD$ & (c) & \ids; independent 60-digit derivatives\\
Correct factor product and normalization & (c) & \edge; 70-digit product\\
Momentum logarithm and $d_0$ & (c) & \edge; symbolic inverse transform\\
Edge energy and analytic $1/(2\pi)$ & (c) & \edge; algebra and two inversion methods\\
Numerical constant and error decomposition & (c) & \ana; free intercepts\\
Uniform limit interchange; full continuum theory & (d) & not established\\
\bottomrule
\end{tabular}
\end{center}

\paragraph{Failures recorded.}
\begin{enumerate}
\item The first scan used $N=\lceil8/h+160\rceil$ and failed the exact
$|P(q)-\pi h|<2\,10^{-10}$ check at $h=.02$. The scan was restarted
with $N=\lceil10/h+200\rceil$. The tolerance was retained, not relaxed.
\item The older half-line note's printed Fourier scale and missing
$\sqrt\pi$ in the momentum transform fail the independent factor-product
and classical-normalization tests. Neither printed expression was used
as an input; the corrected formulas are derived above.
\item A numerical fit cannot establish exact equality of constants, and a
uniform finite-$h$ interchange of limits was not proved. These limitations
are retained despite the successful analytic edge calculation.
\item No historical checker is newly certified. Existing files, README,
and git state are untouched. No final assertion tolerance was weakened.
\end{enumerate}

\paragraph{Reproduction.}
From \texttt{phase2/scripts/}, run
\begin{verbatim}
OPENBLAS_NUM_THREADS=1 OMP_NUM_THREADS=1 \
 PYTHONDONTWRITEBYTECODE=1 python3 onebody_prec_scan.py --recompute
PYTHONDONTWRITEBYTECODE=1 python3 onebody_prec_edge.py --recompute
PYTHONDONTWRITEBYTECODE=1 python3 onebody_prec_identities.py --recompute
PYTHONDONTWRITEBYTECODE=1 python3 onebody_prec_analysis.py
\end{verbatim}
The three numerical scripts validate their saved JSON when \texttt{--recompute}
is omitted. All four terminate with \texttt{ALL CHECKS PASSED}.
Raw data, fit alternatives, algorithm comparisons, source hashes, and
runtimes are retained in \texttt{onebody\_prec\_*.json}.
The recorded dense scan took 1097.7 seconds, the two-method half-line inversions 173.5 seconds, and the high-precision identity check 88.6 seconds; the combined numerical runtime is about 22.7 minutes. Cached validation avoids the dense solves and inverse transforms.
Build with \texttt{pdflatex} twice in \texttt{phase2/}; generated auxiliary
files are kept outside the repository. The delivered source is self-contained.
\end{document}
